% ============================================================
% Gruppenkarten, Geheimkarten, Stationsschilder
% ============================================================
\newcommand{\rollenbox}[1]{%
\begin{metrobox}[title={Rollen in eurer Gruppe},fonttitle=\bfseries]{#1}
\textbf{Buchführung} (1 Person): verwaltet Marker und Bilanzblatt. \textbf{Delegation} (2 bis 3 Personen): verhandelt am Konferenztisch. \textbf{Sprecher\slash in}: vertritt die Gruppe im Plenum und unterschreibt. Alle anderen beraten, entwerfen Angebote und prüfen Verträge. Wechselt die Rollen, wenn ihr wollt.
\end{metrobox}}

% ---------------- Gruppe 1 ----------------
\clearpage
\gruppentitel{cN}{1}{Die Hydroponiker}{Gewächshaus und Pilzkeller}{\iconN{white}{0.9}}
\noindent\textbf{Euer Gut: Nahrung} \rN\par\vspace{1mm}
Ihr betreibt die einzigen noch funktionierenden Anbauflächen der Station. Pilze, Algen und einige genießbare Pflanzen sichern die Ernährung der Bevölkerung. Die Ernte ist begrenzt und wächst nur langsam nach. Ihr seid für alle unverzichtbar, braucht aber selbst Wasser, Energie und Schutzmaterial, um weiter produzieren zu können.
\werte{6 & -- & -- & -- & --}{6 & -- & -- & -- & --}{1 & 2 & 2 & -- & 1}
\textbf{Das braucht ihr von den anderen (je Runde):} 2~\rW\ Wasser, 2~\rE\ Energie, 1~\rS\ Schutzmaterial.

\medskip\noindent\textbf{Eure Interessen}
\begin{itemize}
  \item Niemand darf eure Ernte ohne Gegenleistung beschlagnahmen.
  \item Die Nahrungsmittel müssen gerecht verteilt werden.
  \item Ihr wollt langfristig genügend Vorräte für schlechte Zeiten anlegen.
\end{itemize}
\textbf{Euer Dilemma:} Soll Nahrung nach Bedarf verteilt werden oder sollen diejenigen mehr erhalten, die besonders viel zur Gemeinschaft beitragen?\par
\textbf{Euer Machtmittel:} Ohne eure Ernte droht der Station der Hunger.\par\medskip
\rollenbox{cN}

% ---------------- Gruppe 2 ----------------
\clearpage
\gruppentitel{cW}{2}{Die Wasserwerker}{Pumpenraum und Filteranlage}{\iconW{white}{0.9}}
\noindent\textbf{Euer Gut: Sauberes Wasser} \rW\par\vspace{1mm}
Ihr kontrolliert die letzte funktionsfähige Filteranlage. Ungefiltertes Wasser macht krank. Die Anlage ist beschädigt und braucht Energie und Schutzmaterial (Filter, Masken). Nur ihr kennt alle Arbeitsschritte, um sie sicher zu betreiben.
\werte{-- & 6 & -- & -- & --}{-- & 6 & -- & -- & --}{1 & 1 & 2 & -- & 1}
\textbf{Das braucht ihr von den anderen (je Runde):} 1~\rN\ Nahrung, 2~\rE\ Energie, 1~\rS\ Schutzmaterial.

\medskip\noindent\textbf{Eure Interessen}
\begin{itemize}
  \item Die Anlage muss repariert und gewartet werden.
  \item Niemand darf eigenmächtig Wasserleitungen anzapfen.
  \item Die Wasserversorgung soll nach nachvollziehbaren Regeln erfolgen.
\end{itemize}
\textbf{Euer Dilemma:} Soll jeder Mensch die gleiche Menge Wasser erhalten oder sollen Kranke, Kinder und besonders hart Arbeitende bevorzugt werden?\par
\textbf{Euer Machtmittel:} Ohne euch drohen Durst und Krankheiten.\par\medskip
\rollenbox{cW}

% ---------------- Gruppe 3 ----------------
\clearpage
\gruppentitel{cE}{3}{Die Techniker}{Kraftwerk und Werkstatt}{\iconE{white}{0.9}}
\noindent\textbf{Euer Gut: Energie (Strom und Werkzeug)} \rE\par\vspace{1mm}
Ihr kontrolliert den einzigen noch betriebsfähigen Generator und den größten Bestand an Werkzeugen. Strom ist knapp: Er reicht nicht für Wasseraufbereitung, Beleuchtung, Lüftung und Anbauanlagen gleichzeitig. Auch euer Generator braucht Wartung, Nahrung und Wasser für die Mannschaft sowie Schutzmaterial für Einsätze im Tunnel.
\werte{-- & -- & 5 & -- & --}{-- & -- & 6 & -- & --}{1 & 1 & 1 & -- & 1}
\textbf{Das braucht ihr von den anderen (je Runde):} 1~\rN\ Nahrung, 1~\rW\ Wasser, 1~\rS\ Schutzmaterial.

\medskip\noindent\textbf{Eure Interessen}
\begin{itemize}
  \item Die Energieversorgung muss zentral geplant werden.
  \item Werkzeuge und Ersatzteile dürfen nicht verschwendet werden.
  \item Fachwissen muss geschützt, aber langfristig auch weitergegeben werden.
\end{itemize}
\textbf{Euer Dilemma:} Wer entscheidet, welche Bereiche zuerst Strom bekommen? Fachleute, gewählte Vertreter oder die gesamte Bevölkerung?\par
\textbf{Euer Machtmittel:} Ohne Energie versagen die lebenswichtigen Systeme.\par\medskip
\rollenbox{cE}

% ---------------- Gruppe 4 ----------------
\clearpage
\gruppentitel{cM}{4}{Die Sanitäter}{Krankenstation und Quarantänebereich}{\iconM{white}{0.9}}
\noindent\textbf{Euer Gut: Medikamente und medizinisches Wissen} \rM\par\vspace{1mm}
Ihr verfügt über die letzten größeren Medikamentenvorräte, Verbandsmaterial und Kenntnisse zur Behandlung von Infektionen und Verletzungen. Aus Restbeständen, Pilzkulturen und Heilkräutern stellt ihr in kleinen Mengen Ersatz her. Ohne sauberes Wasser und Strom kann die Krankenstation kaum arbeiten. Jede Runde fordert eine \textbf{Krankheitskarte} Medikamente von der Station.
\werte{-- & -- & -- & 4 & --}{-- & -- & -- & 3 & --}{1 & 1 & 2 & -- & --}
\textbf{Das braucht ihr von den anderen (je Runde):} 1~\rN\ Nahrung, 1~\rW\ Wasser, 2~\rE\ Energie.

\medskip\noindent\textbf{Eure Interessen}
\begin{itemize}
  \item Medizinische Versorgung darf nicht vom persönlichen Reichtum abhängen.
  \item Vorräte müssen für Notfälle zurückgehalten werden.
  \item Medizinische Entscheidungen sollen nicht von politischen Machthabern diktiert werden.
\end{itemize}
\textbf{Euer Dilemma:} Sollen Medikamente zuerst an die Schwerkranken gehen, an die mit den besten Überlebenschancen oder an Menschen, die für die Versorgung der Station besonders wichtig sind?\par
\textbf{Euer Machtmittel:} Ohne euch steigen Krankheit und Sterblichkeit.\par\medskip
\rollenbox{cM}

% ---------------- Gruppe 5 ----------------
\clearpage
\gruppentitel{cS}{5}{Die Wächter}{Zugangsschleusen und Lager}{\iconS{white}{0.9}}
\noindent\textbf{Euer Gut: Sicherheit und Schutzmaterial} \rS\par\vspace{1mm}
Ihr bewacht die Eingänge und kontrolliert die wenigen verbliebenen Vorräte an Schutzkleidung, Filtern und Atemschutzmasken. Ihr sichert gefährliche Bereiche ab und schützt Transporte. Ohne euch könnten Vorräte gestohlen werden oder Gefahren aus den Tunneln eindringen.
\werte{-- & -- & -- & -- & 5}{-- & -- & -- & -- & 3}{1 & 1 & 1 & -- & --}
\textbf{Das braucht ihr von den anderen (je Runde):} 1~\rN\ Nahrung, 1~\rW\ Wasser, 1~\rE\ Energie.

\medskip\noindent\textbf{Eure Interessen}
\begin{itemize}
  \item Die Station braucht verbindliche Regeln und eine durchsetzungsfähige Ordnung.
  \item Die Sicherheitskräfte müssen einsatzfähig bleiben.
  \item Niemand darf die Sicherheit der gesamten Station durch leichtfertige Entscheidungen gefährden.
\end{itemize}
\textbf{Euer Dilemma:} Wie viel Macht darf eine Sicherheitsorganisation besitzen? Soll sie der Regierung unterstehen, unabhängig handeln oder von der Bevölkerung kontrolliert werden?\par
\textbf{Euer Machtmittel:} Ohne euch sind Menschen und Vorräte gefährdet.\par\medskip
\begin{metrobox}[title={Euer Sonderrecht: die Zwangsmaßnahme},fonttitle=\bfseries]{cM}
Einmal pro Spiel dürft ihr bis zu 2 Marker einer Gruppe einziehen. Mit Beschluss der Konferenz ist das legitim. Ohne Beschluss verliert die Station 2~Stabilität.
\end{metrobox}\par\medskip
\rollenbox{cS}

% ---------------- Geheimkarten ----------------
\clearpage
\abschnitt{Geheimkarten \textnormal{\normalsize (ausschneiden, verdeckt austeilen)}}
Jede Gruppe erhält \emph{ihre} Karte. Der Auftrag bleibt geheim, bis die Auswertung beginnt. Erfüllt: \textbf{3 Gruppenpunkte}. Gilt nur, wenn die Station nicht zerfällt.\par\medskip

\begin{schneidbox}[title={Geheimkarte -- nur für Gruppe 1: Hydroponiker},fonttitle=\bfseries]{cN}
\textbf{Eure Ernte ist euer Faustpfand.} Setzt durch, dass im Vertrag steht: Nahrung darf nur gegen Gegenleistung oder mit eurer Zustimmung eingezogen werden. \emph{Rote Linie: eine Regelung, die eure Ernte ohne euch verteilt.}
\end{schneidbox}\par\smallskip
\begin{schneidbox}[title={Geheimkarte -- nur für Gruppe 2: Wasserwerker},fonttitle=\bfseries]{cW}
\textbf{Wasser nach festen Regeln.} Setzt durch, dass die Wasserverteilung nach schriftlichen, öffentlichen Regeln erfolgt und dass ihr in jedem Entscheidungsgremium einen festen Sitz habt. \emph{Rote Linie: Wasser wird zum Druckmittel einer Einzelperson.}
\end{schneidbox}\par\smallskip
\begin{schneidbox}[title={Geheimkarte -- nur für Gruppe 3: Techniker},fonttitle=\bfseries]{cE}
\textbf{Energie bleibt Fachsache.} Setzt durch, dass bei jeder Entscheidung über die Energieverteilung die Techniker mitentscheiden (Sitz oder Vetorecht). \emph{Rote Linie: Energie wird ohne euch verteilt.}
\end{schneidbox}\par\smallskip
\begin{schneidbox}[title={Geheimkarte -- nur für Gruppe 4: Sanitäter},fonttitle=\bfseries]{cM}
\textbf{Medizin kennt keine Parteien.} Setzt durch, dass Medikamente nach medizinischem Bedarf vergeben werden und die Regierung keine Medikamente für einzelne Personen oder Gruppen anordnen darf. \emph{Rote Linie: politische Vorzugsbehandlung.}
\end{schneidbox}\par\smallskip
\begin{schneidbox}[title={Geheimkarte -- nur für Gruppe 5: Wächter},fonttitle=\bfseries]{cS}
\textbf{Ordnung braucht Durchsetzungskraft.} Setzt durch, dass die Wächter im Notstand eigenständig handeln dürfen und nicht von einer einzelnen anderen Gruppe abgesetzt werden können. \emph{Rote Linie: die Wächter werden entwaffnet oder abhängig.}
\end{schneidbox}

% ---------------- Stationsschilder ----------------
\clearpage
\abschnitt{Stationsschilder \textnormal{\normalsize (an die Gruppentische hängen)}}
\foreach \farbe/\name/\ort/\sym in {cN/Hydroponiker/Gewächshaus und Pilzkeller/iconN, cW/Wasserwerker/Pumpenraum und Filteranlage/iconW, cE/Techniker/Kraftwerk und Werkstatt/iconE, cM/Sanitäter/Krankenstation und Quarantäne/iconM, cS/Wächter/Zugangsschleusen und Lager/iconS}{%
\begin{tcolorbox}[enhanced,colback=\farbe!85!black,colframe=black,coltext=white,arc=0mm,boxrule=2pt,height=3.5cm,valign=center,halign=flush center]
{\fontsize{26}{30}\selectfont\bfseries ENDSTATION 7}\\[2mm]
{\Huge\bfseries \name}\hspace{4mm}\csname\sym\endcsname{white}{1.2}\\[1mm]
{\large \ort}
\end{tcolorbox}\par\vspace{1mm}}
